\chapter{系统总体设计}

\section{总体设计思路}

前述需求分析将舵机测试软件的功能归纳为一条连续的测试链路：通信接入提供数据入口，状态确认保证测试条件，参数面板组织不同测试任务，后台过程完成控制和采样，结果数据最终进入显示、保存和复查过程。总体设计需要解决的核心问题，是如何在同一桌面软件中组织上述环节，并在连续反馈数据到达时保持界面响应和测试状态一致。

桌面测控软件通常需要在设备通信、数据采集、界面显示和结果管理之间建立清晰接口\cite{zhang2017qtflowmeter}。数据采集软件相关研究也指出，采集、处理、显示和分析之间的职责边界会影响实时性和可维护性\cite{stansbury2023autodidaqt,turkmanovic2023highspeed}。结合精密舵机测试场景，所设计的软件系统采用分层架构组织总体功能：界面层负责操作输入和结果呈现，桥接层限定界面能够访问的本地能力，协调层统一调度串口通信、测试任务和文件操作，后台任务层承担解析、自动控制和渲染数据整理，串口通信层面向硬件链路提供连接、收发和状态事件。

这种设计的重点在于分离不同频率和不同职责的任务。参数输入、测试项选择和结果查看属于低频人机交互，应主要留在界面层；串口收发、反馈解析、测试任务执行和曲线数据整理具有更强的实时性或计算性，应由后台侧承担。前端通过受控接口获得结构化状态和测试结果，测试核心则通过事件把进度、统计值和完成结果反馈给界面。由此形成的总体结构为通信协议设计和测试算法设计提供了稳定边界，也使后续功能能够围绕同一测试链路展开。

\section{系统总体架构}

系统总体架构可以从数据流方向理解。操作人员在界面层完成串口连接、参数输入、测试启动、取消和数据保存等操作；桥接层将这些操作封装为受控请求；协调层接收请求后调用串口管理、测试核心或后台任务；串口通信层与舵机控制器交换控制帧和反馈帧；处理后的数据再通过事件回到界面层，由示波器、数据显示区域和结果面板呈现。系统总体架构如图~\ref{fig:system-architecture} 所示。

\begin{figure}[htb]
  \centering
  \resizebox{\textwidth}{!}{%
  \begin{tikzpicture}[
    >=latex,
    block/.style={draw, align=center, minimum height=0.85cm, text width=2.45cm, font=\small},
    arrow/.style={->, thick}
  ]
    \node[block] (ui) at (0, 2.4) {界面层\\操作输入与结果呈现};
    \node[block] (bridge) at (3.4, 2.4) {桥接层\\受控接口边界};
    \node[block] (coord) at (6.8, 2.4) {协调层\\任务调度与事件转发};
    \node[block] (worker) at (10.2, 3.7) {后台任务层\\异步处理与数据整理};
    \node[block] (test) at (10.2, 2.4) {测试核心\\测试流程与指标计算};
    \node[block] (serial) at (10.2, 1.1) {串口通信层\\连接收发与帧解析};
    \node[block] (data) at (6.8, 0.2) {结果管理\\保存加载与历史复查};
    \node[block] (device) at (13.6, 1.1) {舵机控制器\\控制帧与反馈帧};

    \draw[arrow] (ui) -- node[above, font=\footnotesize]{操作请求} (bridge);
    \draw[arrow] (bridge) -- node[above, font=\footnotesize]{业务接口} (coord);
    \draw[arrow] (coord) -- (worker);
    \draw[arrow] (coord) -- (test);
    \draw[arrow] (coord) -- (serial);
    \draw[arrow] (serial) -- node[above, font=\footnotesize]{控制帧} (device);
    \draw[arrow] (device) -- node[below, font=\footnotesize]{反馈帧} (serial);
    \draw[arrow] (coord) -- (data);
    \draw[arrow] (data) -- (ui);
    \draw[arrow] (coord) .. controls (6.8, 4.4) and (0, 4.4) .. node[above, font=\footnotesize]{状态、进度与结果事件} (ui);
  \end{tikzpicture}}
  \caption{系统总体架构示意图}
  \label{fig:system-architecture}
\end{figure}

从职责上看，界面层承担通信选择、测试参数输入、控制操作、实时曲线、测试结果显示和数据保存等任务。界面需要同时承载参数录入、状态切换、实时曲线和结果保存等不同粒度的信息，因此采用组件化方式组织通信选择、参数输入、控制、示波器、数据展示和结果保存等区域，各区域围绕同一测试流程协同工作。

桥接层位于界面层和协调层之间，是系统的接口边界。Electron 数据采集桌面应用研究说明，Electron 可以兼顾界面开发效率和本地能力调用\cite{gao2025electron}。在本系统中，本地串口、文件对话框和后台任务不直接暴露给界面层，界面侧只通过桥接层获得有限的业务接口。这样既能满足界面操作需要，也能减少界面层对底层系统能力的直接依赖。

协调层承担系统调度职责。该层向下管理串口连接、控制帧发送、反馈解析、测试任务和文件保存，向上负责把连接状态、解析结果、测试进度和错误信息推送给界面层。协调层的存在使不同测试功能可以共用同一套串口状态、反馈样本和事件回传机制，避免各个面板分别维护硬件访问逻辑。

后台任务层用于承接周期性或耗时处理。当前软件将离线解析、自动控制帧生成和渲染数据整理放入后台任务中执行。后台任务层的设置使高频或周期性工作不直接占用界面线程，同时也让协调层能够以消息方式组织后台处理结果。

串口通信层位于系统与舵机控制器之间。该层负责串口枚举、打开、关闭、写入、连接状态判断、自动重连、控制帧发送和反馈数据解析。串口分包和缓存处理研究说明，连续数据流中的分包、缓存和校验机制会影响通信稳定性\cite{wu2025serialpacket}。因此，本章只说明通信层在总体架构中的位置，具体协议字段、帧格式和解析流程在通信协议与基础控制模块中展开。

\section{进程间通信与事件流设计}

系统内部通信主要通过进程间通信（Inter-Process Communication，IPC）完成。总体上，通信方式可以分为请求响应和事件推送两类。请求响应用于界面层主动发起的操作，例如列出串口、打开串口、启动测试、取消测试和保存数据；事件推送用于协调层主动回传状态，例如串口连接成功、串口断开、数据解析完成、测试进度更新、测试完成和测试错误。IPC 请求与事件流如图~\ref{fig:ipc-event-flow} 所示。

\begin{figure}[htb]
  \centering
  \resizebox{0.96\textwidth}{!}{%
  \begin{tikzpicture}[
    >=latex,
    block/.style={draw, align=center, minimum height=0.82cm, text width=2.55cm, font=\small},
    arrow/.style={->, thick}
  ]
    \node[block] (ui) at (0, 0) {界面操作\\参数与命令};
    \node[block] (bridge) at (3.2, 0) {桥接接口\\请求封装};
    \node[block] (coord) at (6.4, 0) {协调处理\\调度与分发};
    \node[block] (exec) at (9.6, 0.9) {串口与测试任务\\执行请求};
    \node[block] (worker) at (9.6, -0.9) {后台任务\\异步消息};
    \node[block] (event) at (6.4, -2.4) {事件整理\\状态、进度、结果};
    \node[block] (view) at (0, -2.4) {界面更新\\曲线、指标、提示};

    \draw[arrow] (ui) -- node[above, font=\footnotesize]{请求响应} (bridge);
    \draw[arrow] (bridge) -- (coord);
    \draw[arrow] (coord) -- (exec);
    \draw[arrow] (coord) -- (worker);
    \draw[arrow] (exec) -- (event);
    \draw[arrow] (worker) -- (event);
    \draw[arrow] (event) -- node[below, font=\footnotesize]{事件推送} (view);
    \draw[arrow] (view) -- (ui);
  \end{tikzpicture}}
  \caption{IPC 请求与事件流示意图}
  \label{fig:ipc-event-flow}
\end{figure}

表~\ref{tab:ipc-classification} 给出了当前接口的主要分类。表中按业务职责说明接口设计作用，不列出具体通道字符串，以便突出总体设计中的交互边界。

\begin{table}[htb]
  \centering
  \caption{IPC 接口分类及设计作用}
  \label{tab:ipc-classification}
  \small
  \begin{tabular}{@{}p{0.17\textwidth}p{0.14\textwidth}p{0.27\textwidth}p{0.27\textwidth}@{}}
    \toprule
    接口类别 & 主要交互方向 & 承载的操作 & 设计作用 \\
    \midrule
    串口控制接口 & 界面层到协调层 & 串口枚举、连接、断开、控制帧发送 & 为测试流程提供硬件接入入口 \\
    串口状态事件 & 协调层到界面层 & 连接状态、错误状态、解析帧状态 & 让操作人员及时感知通信链路状态 \\
    通信性能测试接口 & 双向 & 测试启动、进度统计、结果回传 & 检查链路质量并决定后续测试条件 \\
    控制性能预测试接口 & 双向 & 噪声估计、阶跃响应、最近结果复用 & 为时域和频域测试提供参数来源 \\
    频域测试接口 & 双向 & 扫频任务、过程结果、频域指标回传 & 支撑频率响应分析和伯德图展示 \\
    时域测试接口 & 双向 & 路径测试、单点统计、时域指标回传 & 支撑位置响应过程和动态指标分析 \\
    数据保存接口 & 界面层到协调层 & 文件选择、结果保存、历史加载 & 将测试结果转化为可追溯数据 \\
    \bottomrule
  \end{tabular}
\end{table}

这种接口分组具有两个方面的作用。其一，界面组件只需要面向所属业务能力发起操作，不必了解协调层内部如何组织串口、测试核心和后台任务。其二，测试核心可以通过统一事件模型向界面反馈状态。以通信性能测试为例，界面发起联合测试请求后，协调层启动通信测试核心；测试核心在运行过程中产生进度、统计、完成或错误事件；协调层再把这些事件转发到界面层。预测试、时域测试和频域测试也采用相近模式，因此不同测试功能在界面层表现出一致的启动、进度、取消和结果展示逻辑。

\section{前端界面结构设计}

前端界面结构围绕一次测试操作的自然顺序组织。主界面划分为左侧操作区、中部示波器区和右下结果管理区。左侧操作区从上到下依次放置通信选择、参数输入和控制操作；中部区域主要承载实时示波器曲线；右下区域由数据显示和数据保存组成。前端界面布局如图~\ref{fig:frontend-layout} 所示。

\begin{figure}[htb]
  \centering
  \includegraphics[width=0.96\textwidth]{chapter3-layout.png}
  \caption{前端界面布局截图}
  \label{fig:frontend-layout}
\end{figure}

通信选择区域是测试流程的起点。该区域负责串口列表加载、端口选择、波特率设置、连接和断开操作，并在连接成功后触发使能检测。前端保留使能检测重试逻辑，并将检测结果同步到共享状态，使数据显示区域能够根据舵机是否使能决定是否显示有效遥测值。这样，连接状态和执行对象状态能够在进入正式测试前被界面明确呈现。

参数输入区域承接通信确认之后的测试任务选择。当前系统按下拉菜单组织四类测试功能，顺序为通信性能测试、控制性能预测试、时域性能测试和频域性能测试。每类测试面板只暴露当前任务所需参数，并通过业务接口调用对应后端能力。面板切换时保留组件状态，可以减少反复切换导致的参数丢失。

控制区域用于处理中止、复位、退出和显示模式等操作。中止操作会按顺序停止自动波形、取消四类可能正在运行的测试任务，并关闭串口连接。这一设计与安全性需求对应，避免自动测试或后台任务在界面退出后继续运行。

数据显示和示波器区域承担过程观察功能。示波器区域接收渲染后的曲线数据，更新四路设定值和反馈值的变化过程。数据显示区域同时接收控制发送信息、实时渲染数据和各类测试完成事件，用于显示实时遥测值和通信、预测试、时域、频域结果摘要。这样，过程数据和结果数据在界面中形成连续展示关系。

数据保存区域承担结果留档功能。该区域从共享结果状态中收集通信测试、预测试、时域测试和频域测试结果，通过数据保存接口调用系统文件对话框和文件写入能力，将结果保存为结构化数据。历史数据加载后会写回共享状态，并触发相关组件恢复显示。图片保存功能在界面中存在入口，但不作为当前论文的已验证核心能力展开。

\section{后端任务调度与 Worker 设计}

后端结构以协调层为中心。协调层注册各类请求处理函数，持有串口管理能力和测试核心，并负责把后台事件转发给界面层。在串口控制链路中，协调层接收打开串口请求后建立串口管理、发送和接收对象；串口收到数据后，协调层组织解析和状态更新，并把原始数据、解析帧或错误信息发送到前端。对于控制帧发送过程，协调层根据当前协议选择对应发送逻辑，并等待发送确认后广播控制值。

测试任务调度同样由协调层统一组织。通信性能测试、控制性能预测试、频域测试和时域测试分别由相应测试核心承担。测试核心通过运行时接口获得发送控制值、读取最近反馈样本、延时和当前时间等能力。协调层把这些能力适配给测试核心，再把测试核心产生的事件转换为界面事件。这样，测试算法可以保持相对独立，界面逻辑也无需介入测试核心内部。

Worker 设计主要服务于后台任务隔离。自动控制运行时，后台任务按指定频率生成控制帧，并通过消息把四路控制值传回协调层；协调层确认当前处于自动控制模式后再发送到串口。离线解析任务可以接收文本或十六进制数据，完成批处理解析和协议帧构造。渲染数据整理任务接收协调层整理后的曲线载荷，并将其转发给界面层。Worker 与协调层之间通过消息通信，可以减少后台处理对界面线程的影响。

这种调度方式使后端形成三类边界：串口管理面向物理链路，测试核心面向业务流程，Worker 面向异步任务。协调层位于这些边界之间，负责把请求、事件、控制帧和反馈样本组织成统一执行过程。通信协议与基础控制模块将在这一框架下进一步展开串口通信和协议处理，测试功能与算法设计则在这一总体调度框架上说明各测试核心的算法流程。

\section{数据流与状态管理设计}

系统的数据流可以分为实时反馈流、测试任务流和结果管理流。实时反馈流从舵机控制器经 RS422 串口进入协调层，经过帧解析后形成反馈数据，再被转换为渲染载荷推送到示波器和数据显示区域。该数据流主要服务于过程观察和异常判断。

测试任务流从前端参数区域发起。用户选择测试类型并填写参数后，前端通过业务接口调用验证和启动能力；协调层启动对应测试核心；测试核心运行过程中通过事件返回进度和统计信息；测试完成后，结果进入前端面板和共享结果状态。通信性能测试、控制性能预测试、时域性能测试和频域性能测试虽然内部算法不同，但在总体数据流上都遵循参数输入、任务启动、事件反馈和结果保存的结构。

结果管理流负责把测试结果从临时界面状态转化为可复查数据。系统为四类测试保存最近一次结果，数据保存区域从共享状态中汇总可用数据，并通过数据保存接口写入本地结构化文件。加载历史数据时，保存区域把文件内容解析后写回共享状态，并通知相关组件刷新。这一设计使实时测试和历史复查可以共用同一套结果展示逻辑。

三类数据流之间存在递进关系。实时反馈流保证测试过程可观测，测试任务流将反馈样本转化为可解释指标，结果管理流则把测试输出保存为可追溯对象。总体设计通过 IPC、共享状态和 Worker 消息把这些流连接起来，使界面操作、后台执行和结果留档保持一致。

\section{本章小结}

本章在需求分析的基础上给出了舵机测试软件的总体设计。系统采用以界面层、桥接层、协调层、后台任务层、串口通信层和测试核心为主体的分层架构，将界面交互、本地硬件访问、后台任务、协议通信和测试流程分离开来。这样的结构能够支撑连接管理、状态确认、参数输入、四类测试、实时显示和结果保存之间的连续衔接。

在通信机制上，系统通过请求响应和事件推送两类 IPC 形式组织前后端协作。前端界面围绕测试操作顺序布置，后端协调层负责串口、测试核心和 Worker 的统一调度，数据流则从实时反馈、测试任务和结果管理三个方向组织。上述总体设计为后续通信协议与基础控制模块设计提供了硬件链路和数据边界，也为测试功能与算法设计提供了任务调度和结果回传框架。
